\subsection{Gravitational Wave Detection and Modern Digital Constraints}
The direct observation of gravitational waves fundamentally relies on advanced signal processing to isolate transient astrophysical phenomena from environmental noise \cite{1}. In the primary literature, this extraction is predominantly executed using complex digital signal processing (DSP). These systems utilize matched filtering algorithms and adaptive digital notch filters to dynamically track and eliminate 60~Hz power-line harmonics and seismic interference \cite{1}. While effective, these digital methods impose a substantial computational demand and operate almost entirely in the discrete-time domain post-digitization.

The prevalence of these digital back-ends is supported by the standardized access provided by the LIGO Open Science Center (LOSC), which facilitates the deployment of coherent search algorithms within non-Gaussian noise floors \cite{6, 10}. However, because these methodologies depend on post-digitization processing, they are inherently constrained by the dynamic range of the initial sampling hardware. This architectural landscape suggests that while digital identification remains the standard for statistical verification, it is primarily reactive to terrestrial interference that has already permeated the data stream.

\subsection{Physical Modeling and the LTI Framework}
The theoretical foundation for front-end, continuous-time signal conditioning relies on classical systems engineering literature. Standard texts define the physical constraints of Linear Time-Invariant (LTI) systems and establish the mathematical principles of convolution used to model environmental distortion \cite{2}. Under the LTI framework, the signal path from the gravitational event through the interferometric sensors operates as a linear system where the distorted output $y(t)$ is the convolution of the pure astrophysical signal $x(t)$ with the system’s impulse response $h(t)$, represented by $y(t) = x(t) * h(t)$.

This modeling approach appears particularly relevant for binary black hole mergers, which astrophysical theory characterizes as a linear frequency-modulated (LFM) chirp where orbital decay dictates a monotonic increase in frequency and amplitude \cite{7}. By treating the observatory as a deterministic LTI system \cite{2}, the environmental distortion of this known chirp waveform \cite{7} is viewed as a physical convolution. This theoretical alignment justifies the application of inverse filtering, which seeks to mathematically reverse the environmental influence before the signal is processed by stochastic digital identification algorithms.

\subsection{Foundations of Filter Topologies and Design Trade-offs}
Established analog filter design literature outlines the fundamental mathematical trade-offs required to condition continuous signals. These texts strictly define the relationships between magnitude response, phase distortion, and roll-off steepness inherent in classical Butterworth, Chebyshev, and Elliptic filter topologies \cite{3}. Standard design principles categorize Butterworth filters by their maximally flat passband, preserving target frequencies without amplitude distortion at the cost of a gradual transition band \cite{9}.

Conversely, Chebyshev and Elliptic architectures introduce equiripple behavior to achieve highly aggressive transition characteristics \cite{3, 9}. Technical consensus in analog theory indicates that no single topology is universally superior; the architecture must be selected based on the spectral proximity of the interference. Given the significant amplitude of the 60~Hz interference adjacent to the 30--300~Hz astrophysical band, the sharp magnitude roll-off of an Elliptic topology emerges as a necessary compromise. It achieves the required attenuation, even if it introduces non-linearities that simpler designs omit.

\subsection{Deconvolution and the Challenges of Inverse Filtering}
Restoring the original signal trajectory requires the mathematical inversion of the environmental convolution through spectral deconvolution. According to the principles of discrete-time signal processing, this involves dividing the Fourier transform of the captured data by the system’s frequency response: $\hat{X}(f) = Y(f)/H(f)$ \cite{4}. While a standard tool for signal restoration, it remains an ill-posed problem highly susceptible to noise enhancement \cite{4}.

In frequency regions where the detector's sensitivity $H(f)$ is low, division by small values causes background noise to amplify exponentially \cite{4, 5}. This phenomenon is particularly problematic for LFM chirps, as their broadband nature forces them across these low-sensitivity spectral regions. Consequently, the recovery process cannot rely on deconvolution alone; it necessitates high-order filtering to precisely excise the high-frequency artifacts generated by the inversion process.

\subsection{Discrete Realization via Bilinear Mapping}
To bridge the gap between continuous-time analog theory and modern digital simulation, literature provides mapping techniques such as the Bilinear Transformation. This numerical method maps the s-plane into the unit circle of the z-plane, ensuring that the stability of a physical analog filter is preserved within a digital environment \cite{5}. A characteristic of this mapping is frequency warping, a non-linear relationship between analog and digital frequencies that increases as the signal approaches the Nyquist limit \cite{5, 8}.

Standard DSP texts suggest that while warping must be accounted for via pre-warping techniques, the Bilinear Transformation remains a highly reliable method for modeling physical hardware responses \cite{4, 8}. For the low-frequency 30--300~Hz band utilized in black hole detection, the warping effect is minimal. This allows researchers to evaluate the efficacy of continuous-time architectures \cite{3} within a controlled digital simulation environment without losing the underlying physical relevance of the analog design.

\subsection{Synthesis}
This project distinguishes itself from primary modern works by focusing strictly on the capabilities of foundational analog filtering. While prior work frequently relies on adaptive digital algorithms to retroactively cancel 60~Hz electromagnetic interference \cite{1}, this study restricts the solution to front-end continuous-time architecture. It investigates whether static, continuous-time analog bandpass geometries can independently achieve the attenuation necessary to isolate a transient chirp. By applying classical analog design principles \cite{3} directly to empirical strain data \cite{1}, this research isolates the exact mathematical filter tolerances specifically the passband and stopband ripples of the Elliptic filter required to physically separate a cosmic signal from severe electromagnetic interference prior to advanced digital processing.

The technical literature reveals a fundamental tension between computationally intensive digital identification strategies \cite{1, 10} and the deterministic potential of foundational systems theory \cite{2}. By unifying the mathematical precision of analog filter topologies \cite{3, 9} with the practical realization enabled by Bilinear mapping \cite{5}, this study evaluates whether the physical constraints of the detector can be pre-emptively addressed. This approach shifts the detection paradigm from a reactive digital strategy toward an architectural front-end method, offering a mathematically rigorous framework for signal isolation in high-interference environments.